% Deutsche Parallelfassung zu state/state_taxonomy.tex

\chapter{Zustandsbegriffe als qualifizierte Verwendungen}
\label{chap:state-taxonomy}
\index{Zustand!qualifizierte Verwendung}

Die Thesis verwendet das Wort \emph{Zustand} für mehrere Ausgaben und
Sichtweisen. Dieses Kapitel ist weder eine zweite Ontologie noch eine
kanonische Hierarchie. Es ist eine Leselandkarte: Jeder folgende Begriff
übernimmt die in den Teilen III--V eingeführten Qualifikationen zu Objekt,
Zeit, Realisierung, Provenienz, Anwendbarkeit und Zweck.

\begin{principle}[Zustandsschichtprinzip]
\label{princ:state-layer-principle}
Ein Zustandsbegriff ist nur zusammen mit Engineering-Frage, Gegenstand,
Bedingungen und Gültigkeit aussagekräftig. Die gemeinsame Verwendung des
Wortes ``Zustand'' macht verschiedene Ausgaben nicht austauschbar.
\end{principle}

\section{Bereits eingeführte Begriffe}

\paragraph{Zustand.}
\label{def:taxonomy-general-state}
Hier wird keine neue allgemeine Definition eingeführt. In diesem Teil
bezeichnet Zustand eine qualifizierte Antwort zu einem identifizierten
Gegenstand unter angegebenen Bedingungen.

\paragraph{pose2.}
\label{def:taxonomy-pose2}
Es gilt der ebene Zustand aus Definition~\ref{def:pose2}: Er wird aus dem geordneten
Krümmungsgesetz des Alignments und einer Anfangspose unter einer
Realisierungskonvention rekonstruiert. Er ist weder Alignment-Identität noch
ein dauerhafter kanonischer Kernel.

\paragraph{pose3.}
\label{def:taxonomy-pose3}
Es gilt die qualifizierte räumliche Eisenbahnpose aus
Definition~\ref{def:runtime-pose3}. Sie ergänzt Profil, Überhöhung, metrische
Realisierung und Frame-Interpretation, ohne zur vollständigen physischen
Anlage zu werden.

\paragraph{Laufzeitzustand.}
\label{def:taxonomy-runtime-state}
Diese Formulierung bezeichnet ein bei Bedarf ausgewertetes anwendbares
Ergebnis. Sie identifiziert keinen universellen Zustandstyp und verbirgt den
erzeugenden Operator nicht.

\paragraph{Realisierter Zustand.}
\label{def:taxonomy-realized-state}
Der Begriff wird nur für einen durch eine angegebene Realisierung
qualifizierten Zustand verwendet. Entwurfsrealisierung, gebauter physischer
Zustand und Schätzung des beobachteten Zustands bleiben getrennt.

\section{Durch Verwendung oder Modell eingeführte Begriffe}

\paragraph{Betriebszustand.}
\label{def:taxonomy-operational-state}
Eine für eine betriebliche Frage verwendete Ausgabe; die Bezeichnung begründet
für sich weder Komfort, Sicherheit, Zulässigkeit noch Autorität.

\paragraph{Fahrzeugzustand.}
\label{def:taxonomy-vehicle-state}
Ausgabe eines identifizierten Fahrzeugmodells relativ zum qualifizierten
Eisenbahnzustand und zur Durchfahrt, wie in Kapitel~\ref{chap:vehicle-space}
erläutert.

\paragraph{Schwerpunktzustand.}
\label{def:taxonomy-cg-state}
Modellausgabe für einen gewählten Massenbezugspunkt. Sie ist nicht der
vollständige Fahrzeugzustand und benötigt das relevante Massen- und
Fahrzeugmodell.

\paragraph{Dynamischer Zustand.}
\label{def:taxonomy-dynamic-state}
Die in Definition~\ref{def:runtime-dynamic-state} eingeführte modellspezifische
Ergebnisbahn \(x_M\); sie ist kein Synonym für pose3 oder physischen Zustand.

\paragraph{Optimierungszustand.}
\label{def:taxonomy-optimization-state}
Eine bewusst für eine Optimierungsformulierung gewählte Variable oder Ausgabe.
Ihre Auswahl ist eine explizite Engineering-Wahl, und ein Optimum ist keine
Entscheidung.

\paragraph{Relativer Zustand.}
\label{def:taxonomy-relative-state}
Ein relativ zu einem angegebenen Referenzzustand oder Frame ausgedrücktes
Ergebnis. Referenz und Konvention bleiben Teil seiner Bedeutung.

\section{Abhängigkeit statt Hierarchie}

Es gibt keine universelle Kette
\(\kappa\rightarrow\mathrm{pose2}\rightarrow\mathrm{pose3}\rightarrow
x_{\mathrm{vehicle}}\rightarrow x_{\mathrm{opt}}\).
pose2 und pose3 folgen den räumlichen Abhängigkeiten aus Teil V. Durchfahrt
ergänzt Zeit. Ein gewähltes Antwortmodell kann Fahrzeugausgaben ergänzen. Eine
Optimierungsformulierung kann danach eine qualifizierte Größe als Variable,
Ziel oder Nebenbedingung wählen. Andere Verwendungen enden früher.
Abhängigkeit ist verwendungsspezifisch und keine kanonische Zustandsleiter.

\begin{summary}
Diese Landkarte verhindert, dass gleichlautende Begriffe zu
Identitätsbehauptungen werden. Sie führt kein neues Zustandsmodell ein: Jeder
nachgelagerte Begriff bezeichnet eine qualifizierte Verwendung oder
Modellausgabe aus der bereits eingeführten Eisenbahnzustandsdarstellung.
\end{summary}
